\section{타당성 위협과 한계}
\label{sec:threats}

\subsection{구성 타당성}
CoC는 인간이 직접 작성한 안전 명세가 아니라 기록 궤적에 조건화된 VLM/LLM의 사후 주석이다. 따라서 문장이 장면의 원인이나 운전 의도를 충실하게 설명한다는 보장은 없다. \texttt{quality\_pass}, 파서와 의도 매핑은 검증 입력을 구성하는 기계 생성 신호일 뿐 인간 정답을 대체하지 않으며, 본 결과는 일반적인 CoC 충실도가 아니라 구성된 시간ㆍ출처 계약의 내부 정합성만 다룬다.

\subsection{측정 타당성}
응답은 액추에이터 명령이 아니라 자차 궤적에서 계산한 속도 변화 대리변수이다. 창 $W$, 속도 임계값 $\delta$, 방향 비율 $\rho$에 따라 판정이 달라질 수 있고, 27개 설정도 모든 유효 탐지기를 포괄하지 않는다. $v\leq0.5$~m/s STOP/YIELD 이미 충족 규칙은 초기 검토 대상 8건 중 3건을 재분류한 분석 대상 집합 보정 탐색 논리이다. 또한 $D=3$초는 실행 가능성 분석 설정이며 규범적 안전 기한이 아니다. 기한은 $r<D$로 해석하지만 저장 집계에는 정확히 $r=D$인 사례가 없고, 등호 시점 전환 우선순위는 후속 검사기 회귀 시험이 필요하다. \preserve/\reset 반복 정책과 기한 선택에 따른 판정 반전 때문에 계약 통과와 검토 후보는 측정ㆍ정책 조건을 함께 명시해야 한다.

\subsection{검증 조건의 타당성}
의도별 응답 술어, 잠정 제한 시간, 반복 사건 및 종료 시점 의미론은 연구자가 선택한 검증 조건이다. 이는 법규, 제어기 요구사항 또는 차량 동역학에서 독립적으로 도출한 규범 계약이 아니다. STOP/HOLD/RECOVER 다단계의 일부 전환도 합성 가정이므로 관찰자 표현력 시험에는 쓸 수 있지만 실제 행동 의무나 학습 정답으로 해석할 수 없다.

\subsection{내적 타당성}
일괄 결과는 파서, 품질 관문, 의도 매핑, 응답 탐지기와 결과 분류기의 구현에 의존한다. 282개 분석 대상 집합 모델은 위반 불리언을 직접 설정하고 1,128개 기대 판정도 레이블에서 도출하므로 그 일치는 플래그/질의 직렬화 회귀만 지지한다. 관찰자 수준 시간 변이 근거는 \texttt{scene-0001} 모델에 한정하며 미관측 결함의 검출 성능은 측정하지 않는다. 특히 \texttt{orphan\_response}는 합성 연결 오류 조건이며 자연 발생 오류의 관측치가 아니다.

\subsection{외적 타당성}
보조 단일사건 분석은 chunk~0의 94개 장면과 403개 주석 중 Qwen/Qwen3.5-35B-A3B의 자기평가를 통과하고 영어 정규식/복합 의도 규칙에서 방향이 명시된 134개 주석 사건별 검증 조건에 한정된다. 품질 통과 290개 중 모호ㆍ비방향ㆍ복합 의도 156개는 제외했으므로 이 결과는 해당 부분집합에만 적용된다. 연속 CoC 자연 검토도 309개 인접 창 전체가 아니라 자동 후보 15개와 대응 대조 12개로 선택한 27개 장면만 포함하며 무작위 표본이 아니다. 다른 nuScenes 분할, CoC 생성기, 언어, 횡방향 의도, 센서ㆍ차량 및 학습 정책으로 일반화하려면 별도 평가가 필요하다.

\subsection{검토 타당성}
기존 다섯 후보의 전방 카메라 프레임 검토는 독립적인 인간 평가가 아니라 비인간 AI 시각 선별이었다. 연속 CoC 실험에서는 후보ㆍ대조와 검사기 출력을 가린 텍스트 전용 화면에서 네 검토자가 27장면을 독립 판정하고 별도 합의를 수행했다. 그러나 27개 모두 적어도 한 항목에서 의견이 갈렸고 텍스트 일관성의 Fleiss $\kappa$는 0.027로 낮았다. 합의 CSV에는 항목별 판단 근거가 기록되지 않아 합의 과정을 제3자가 재구성할 수도 없다. 따라서 27/27 일관 합의는 선택 표본의 최종 패널 판정일 뿐 독립적인 정답 집합이 아니다. 또한 영상과 객체 상태를 제공하지 않았으므로 합의는 텍스트 내부 요구 관계만 다루며 실제 해제 상태나 행동 적절성의 정답이 아니다.

\subsection{합성 회귀 시험의 타당성}
40쌍은 자연 문장을 변형한 자료가 아니라 코드가 직접 만든 합성 중간표현이다. 생성기는 기준 사례가 상태형 검사에서 모순이 아니고 변형 사례가 지정 모순으로 판정되는지를 사후조건으로 강제한다. 따라서 상태형 40/40과 기준 사례 0/40은 독립 검출 성능이 아니라 생성 규칙과 검사 규칙의 자기 일관성 및 회귀 fixture 커버리지를 나타낸다. 사건별 20/40과의 차이도 과거 상태를 보지 않도록 정의한 기준선의 구조에서 예상되는 결과이다. 자연 발생 모순ㆍ비모순 사례에 대한 정밀도, 재현율 또는 오검출률은 측정하지 않았다.

\subsection{출처 타당성}
로컬 CoC-Nusc 파일은 장면 간 로그ㆍ이전/다음 관계와 객체 추적 연속성을 제공하지 않으므로 기본적으로 독립 에피소드로 닫았다. 13체인 기본 동작 산출물은 공식 trainval 원본의 바이트 수준 출처를 재검증하지 못한 nuScenes 호환 메타데이터 미러에 의존한다. 더구나 느린 의도 우선 부분문자열 파서와 $W=1.0$초, $\delta=0.7$~m/s, 방향 비율 없는 탐지기는 공유 계약과 다르고 복합 의도를 오분류할 수 있다. 따라서 4/9 및 8/7/0은 해당 구현의 기술 출력일 뿐 공식 데이터, 객체ㆍ위험 연속성 또는 장기 계약 성능의 정량 근거가 아니다.

\subsection{정형화 한계}
현재 모델은 기록된 결정적 사건열과 자차 궤적을 재생하며 환경이나 제어 대안을 탐색하지 않는다. 기존 집계 산출물의 13/13 기록과 달리 연속 CoC 실험은 실행 가능한 Python 기준선과 실제 \texttt{verifyta}의 일곱 정식 시험 결과를 보존했다. 그러나 UPPAAL 구현도 같은 전이 규칙에서 생성되므로 0/7 불일치는 제한된 구현 일치 시험이지 의미론의 독립 검증이 아니다. 특히 중심 연속 CoC 관찰자는 committed 전환으로 사건 순서만 소비하고 실제 타임스탬프 차이나 기한을 제약하지 않는다. 계약 의미론 자체의 완전성이나 다른 종류의 자연 오류 검출을 입증하지 않으며, 상대 객체 상태, 액추에이터, 차량 동역학 및 폐루프 상호작용도 포함하지 않으므로 물리 안전성, 제어 합성 또는 배포 적합성을 증명하지 않는다.
